\uuid{iW9n}
\titre{Limite}
\auteur{bodin ; hanani ; mzari}
\organisation{exo7}
\licence{CC BY-NC-SA 4.0}
\datecreate{2020-01-20}
\niveau{L1}
\module{Analyse}
\chapitre{Suite}
\sousChapitre{Autre}
\difficulte{1}
\theme{Suites, limite}
\exo7qcmid{lille-1.suites.006}

\contenu{
\begin{question}
\qclassification{121.00}
\qtags{motcle=limite}
Soit $\displaystyle u_n=n\ln\left(1+\frac{1}{n}\right)$ et $\displaystyle v_n=\left(1+\frac{1}{n}\right)^n$. Quelles sont les bonnes réponses ?
\begin{answers}
    \bad{$\displaystyle \lim _{n\to +\infty}u_n=+\infty$ et $\displaystyle \lim _{n\to +\infty}v_n=1$}
    \bad{$\displaystyle \lim _{n\to +\infty}u_n=0$ et $\displaystyle \lim _{n\to +\infty}v_n=1$}
    \bad{$\displaystyle \lim _{n\to +\infty}u_n=1$ et $\displaystyle \lim _{n\to +\infty}v_n=1$}
    \good{$\displaystyle \lim _{n\to +\infty}u_n=1$ et $\displaystyle \lim _{n\to +\infty}v_n=\mathrm{e}$}
\end{answers}
\begin{explanations}
On utilise le fait que si $\displaystyle \lim _{n\to +\infty}a_n=0$, alors les suites $(a_n)$ et $\ln (1+a_n)$ sont équivalentes. Ainsi le terme $u_n$ est équivalent, en $+\infty$, à $\displaystyle n\times \frac{1}{n}=1$. Donc $\displaystyle \lim _{n\to +\infty}u_n=1$ et, puisque $v_n=\mathrm{e}^{u_n}$, $\displaystyle \lim _{n\to +\infty}v_n=\mathrm{e}$.
\end{explanations}
\end{question}
}
